\scanpage{26}
un sens les foncteurs ``ombres'' $\omega_N$ sont d'autant plus grossiers que $N$ est plus grand. Néanmoins ils sont tous pleinement fidèles (\uline{NB} les foncteurs $\tau_N$ ne sont pas pl. fidèles bien sûr, pas $=$ sur la sous-catégorie des ``ombres virtuelles d'ordre $N-1$'', i.e. des modules gradués dont les composantes sont nulles en degrés $\leq N-1$.

Soit $\operatorname{Omb}(k)$ la catégorie des ``ombres virtuelles'' des $k$-modules, i.e. des $S=k\{T\}$-modules gradués [à degrés $\geq0$?]. Soit, $\forall N\geq0$, $\operatorname{Omb}_{\unsure{N\text{-n}}}(k)$ la sous-catégorie pleine des modules $N$-négligeables, i.e. tels que $\mathcal M_i=0$ si $i\geq N$.
C'est une sous-catégorie épaisse, la catégorie quotient est can. équivalente \deleted{à celle} \inserted{\unsure{la s-cat. pleine de (Omb) formée}} des ``\inserted{$N$-}ombres virtuelles'', i.e. des $\unsure{S|k}$-modules gradués tels que $\mathcal M_i=0$ si $i<N$.
Le foncteur passage au quotient s'identifie alors au foncteur ``troncature'' $\tau_N$. On note que
\diagram{\begin{tikzcd}
\tau_N(\mathcal M) \arrow[r,hook] & \mathcal M
\end{tikzcd}}
(inclusion fonctorielle) \inserted{correspond au fait que} le foncteur \deleted{$\unsure{\tau_N}$} pass. au quotient admet un adjoint à droite (pl. fid. comme de juste)
\[
\operatorname{Hom}_{\operatorname{Omb}}(\mathcal N,\tau_N(\mathcal M))\simeq
\operatorname{Hom}_{\substack{N\text{-Omb}\\\text{ou Omb}}}(\tau_N(\unsure{\mathcal M|\mathcal N}),\tau_N(\mathcal M)).
\]

Appelons ``négligeable'' une ombre $\mathcal M$ qui est $\varinjlim$ \inserted{filtr.} d'ombres $N$-négligeables (pour $N$ variable), i.e. telle que \inserted{$\forall i\in\mathbb Z$,} $\forall x\in\mathcal M_i$, $\exists N\unsure{(\in\mathbb N)}$ tel que $T^{(\alpha)}.x=0$ pour $\alpha\geq N$ i.e. $S^{+(N)}.x=0$. C'est encore une sous-catégorie épaisse, et le localisé est
\mspage{25}
